\documentclass[11pt,a4paper]{article}
\usepackage[utf8]{inputenc}
\usepackage[T1]{fontenc}
\usepackage[english]{babel}
\usepackage{lmodern}
\usepackage{geometry}
\usepackage{xcolor}
\usepackage{tocloft}
\usepackage{hyperref}
\usepackage{enumitem}
\usepackage{parskip}
\usepackage{booktabs}
\usepackage{longtable}
\usepackage{fancyhdr}
\usepackage{lastpage}
\usepackage{listings}
\usepackage{titlesec}
\geometry{margin=2.15cm}
\setlength{\headheight}{14pt}
\setlength{\emergencystretch}{3em}
\setlist{itemsep=3pt,parsep=1pt,topsep=4pt}
\definecolor{titleblue}{HTML}{123B5D}
\definecolor{softblue}{HTML}{EDF5FB}
\definecolor{softgray}{HTML}{F5F7FA}
\definecolor{accent}{HTML}{0F6E8C}
\hypersetup{
  colorlinks=true,linkcolor=titleblue,urlcolor=accent,
  pdftitle={Scientific Data Analysis v1.9: English publication edition},
  pdfauthor={interemi},
  pdfsubject={Contract freeze and app-facing hardening}
}
\titleformat{\section}{\Large\bfseries\color{titleblue}}{\thesection.}{0.6em}{}
\titleformat{\subsection}{\large\bfseries\color{accent}}{\thesubsection.}{0.6em}{}
\setlength{\cftsecnumwidth}{2.8em}
\tocloftpagestyle{fancy}
\lstdefinestyle{promptstyle}{
  basicstyle=\ttfamily\small,backgroundcolor=\color{softgray},
  frame=single,rulecolor=\color{softblue},breaklines=true,columns=fullflexible
}
\pagestyle{fancy}
\fancyhf{}
\fancyhead[L]{Scientific Data Analysis}
\fancyhead[R]{v1.9 guide: English edition}
\fancyfoot[C]{\thepage\ of \pageref{LastPage}}
\newcommand{\code}[1]{\texttt{\detokenize{#1}}}

\begin{document}
\pagenumbering{gobble}
\begin{titlepage}
\centering
\vspace*{1.0cm}
{\Huge\bfseries\color{titleblue} Consolidation Guide\\[0.15cm]
Scientific Data Analysis v1.9\par}
\vspace{0.65cm}
{\Large Contract freeze and app-facing hardening\par}
\vspace{0.8cm}
{\large Stable artifact types, clean errors, P1/P2 consolidation,
and expanded app-like regressions\par}
\vspace{0.8cm}
{\large English publication edition: September 23, 2026\par}
\vspace{0.5cm}
\begin{minipage}{0.93\textwidth}
\small\raggedright
Translation of the preserved private v1.9 guide at commit \code{b9f6176}.
Original TeX SHA-256:
\nolinkurl{37247e3ded372906cdca842b400c81b1264a6ecb343226dac5b50a6be8022ae2}.
Private source:
\nolinkurl{skills/scientific-data-maintainer/.cache/active_references/guides/guia_scientific_data_analysis_v1_9.tex}.

Recorded results and closure steps are historical, not current verification.
This release guide does not synchronize or modify Scientific Workbench;
the app remains outside its scope. Historical commands are retained for
traceability: use fresh output directories when reproducing them, without
overwriting earlier evidence. The original remains unchanged privately.
\end{minipage}
\vfill
{\large Prepared with Codex\par}
{\large Copyright 2026 interemi\par}
\end{titlepage}

\clearpage
\pagenumbering{roman}
\fancyfoot[C]{\thepage}
\begin{abstract}
\code{scientific-data-analysis} v1.9 is a consolidation release. v1.8 prepared
the backend for app consumption; v1.9 reduces debt and freezes the components
an app needs to remain stable: artifact type names, app-facing errors, exposure
limits, app hints, next actions, and realistic regressions.

v1.9 makes the backend more predictable, its blocks clearer, and its outputs
easier to parse before major v2.0 integration. It does not develop or
synchronize Scientific Workbench.
\end{abstract}
\tableofcontents
\newpage
\pagenumbering{arabic}
\fancyfoot[C]{\thepage\ of \pageref{LastPage}}

\section{Executive overview}
v1.9 is the Contract Freeze and App-Facing Hardening phase. It consolidates
the v1.8 work by:
\begin{itemize}[leftmargin=1.5em]
\item Classifying inherited P1/P2 debt as v1.9 work, v2.0 deferrals, or not applicable.
\item Freezing \code{artifact_type} names while retaining published v1.8 compatibility.
\item Hardening errors, warnings, and \code{BLOCKED_CONTROLADO} so apps need not interpret tracebacks.
\item Reducing risky duplication in envelope, artifact typing, app hint, and next-action helpers.
\item Fixing normal, expert, optional, legacy, maintainer, and unexposed modes.
\item Expanding app-like regressions with realistic, ambiguous, and broken cases.
\item Updating the guide and references to explain consolidation rather than claim v2.0.
\end{itemize}
The product decision is to prepare the backend for a controlled future
integration, without creating Scientific Workbench v2.0 in this release.

\section{From v1.8 to v1.9}
\begin{longtable}{p{0.23\linewidth}p{0.65\linewidth}}
\toprule
\textbf{Area} & \textbf{v1.9 change} \\
\midrule
Charter & Defines a contract-freeze and app-facing-hardening release. \\
Debt inventory & Classifies all 59 public capabilities by family, decision, phase, risk, and focused action. \\
Artifact types & Freezes official names and allowed aliases; prohibits generic types such as \code{json}, \code{csv}, \code{png}, and \code{report}. \\
Errors & Requires \code{errors.kind}, \code{errors.message}, \code{next_actions}, \code{original_modified}, and no raw traceback. \\
App hints & Priority routes emit or document \code{short_summary}, \code{severity}, \code{preview_artifact_types}, \code{tags}, and useful actions. \\
Exposure & Optional, legacy, expert, and maintainer routes stay outside normal mode; the router rejects or demotes them. \\
Regressions & Adds professional data, documents, inherited notebooks, corrupt containers, time series, sensors, expert astronomy, missing backends, and output conflicts. \\
\bottomrule
\end{longtable}

\section{What remains unchanged}
\begin{itemize}[leftmargin=1.5em]
\item Scientific Workbench is not edited in v1.9.
\item Installed synchronization waits until the final installation phase.
\item Published v1.8 artifact types are not renamed without compatibility.
\item Domain-specific, legacy, and maintainer tools are not forced into general use.
\item UI-driven workflows, job history, preview panels, and a formal app/skill planner remain outside this release.
\item Optional backends do not become core dependencies.
\item Correct scientific behavior is preserved when improving app-facing output.
\end{itemize}

\section{v1.9 source documents}
\begin{itemize}[leftmargin=1.5em]
\item \path{references/v1-9-consolidation-charter.md}: goals, non-goals, phases, and closure.
\item \path{references/v1-9-debt-inventory.md}: actionable P1/P2 inventory for 59 capabilities.
\item \path{references/v1-9-artifact-types-freeze.md}: official artifact types, aliases, and prohibited names.
\item \path{references/v1-9-error-contract.md}: app-facing errors, failures, and blocks.
\item \path{references/v1-9-app-hints-next-actions.md}: priority-route hints and next actions.
\item \path{references/v1-9-exposure-modes.md}: normal, expert, optional, legacy, maintainer, and unexposed modes.
\item \path{references/v1-9-v2-0-deferrals.md}: protected v2.0 rows.
\item \path{references/v1-9-app-like-extended-regression.md}: expanded app-like coverage.
\end{itemize}

\section{P1/P2 inventory}
The inventory covers 59 public capabilities plus two global rows. It classifies
app-facing debt as closed, deferred, or explicitly inapplicable; it does not
mean that the skill is broken.

\begin{longtable}{p{0.20\linewidth}r p{0.59\linewidth}}
\toprule
\textbf{Severity} & \textbf{Rows} & \textbf{Interpretation} \\
\midrule
\code{P1} & 25 & Significant app-facing, optional, artifact-typing, or future UX work. \\
\code{P2} & 34 & Consolidation, documentation, tests, hints, limits, or regression improvements. \\
\code{P0} & 0 & No critical release blocker documented in the inventory. \\
\bottomrule
\end{longtable}

\begin{longtable}{p{0.20\linewidth}r p{0.59\linewidth}}
\toprule
\textbf{Decision} & \textbf{Rows} & \textbf{Meaning} \\
\midrule
\code{v1_9} & 34 & Focused consolidation changes. \\
\code{v2_0} & 8 & Requires Scientific Workbench, UI, or substantial workflow work. \\
\code{no_procede} & 17 & Close by documenting limits and accurate exposure. \\
\bottomrule
\end{longtable}

\subsection{Debt families}
\begin{longtable}{p{0.28\linewidth}p{0.60\linewidth}}
\toprule
\textbf{Family} & \textbf{Consolidation scope} \\
\midrule
\code{artifact_types} & Stable typing for outputs, previews, manifests, reports, visual FITS, and sidecars. \\
\code{app_hints} & Short summaries, severity, tags, previews, and useful next actions. \\
\code{errors} & Clean, parseable failures/blocks without raw tracebacks. \\
\code{exposure_modes} & Normal, expert, optional, legacy, and maintainer visibility. \\
\code{duplicacion} & Duplicate helpers affecting envelopes, artifacts, manifests, or next actions. \\
\code{regressions} & Expanded, lightweight, deterministic app-like coverage. \\
\bottomrule
\end{longtable}

\section{Frozen artifact types}
Primary reference: \path{references/v1-9-artifact-types-freeze.md}.

\subsection{Inherited v1.8 types}
\begin{longtable}{p{0.28\linewidth}p{0.60\linewidth}}
\toprule
\textbf{Artifact type} & \textbf{Use} \\
\midrule
\code{summary_json} & Parseable run summary or envelope. \\
\code{manifest_json} & Provenance, input, output, and environment manifest. \\
\code{report_md} & Human-readable Markdown report. \\
\code{preview_png} & Static visual or Quick Look preview. \\
\code{table_csv} & Derived CSV or TSV table. \\
\code{notebook_ipynb} & Generated, inspected, or copy-executed notebook. \\
\code{log_txt} & stdout, stderr, command transcript, or logs. \\
\code{qa_report} & Dedicated QA report. \\
\code{handoff_bundle} & Folder/package for delivery or continuation. \\
\code{unknown} & Explicit fallback when better typing is unavailable. \\
\bottomrule
\end{longtable}

\subsection{Compatible v1.9 additions}
\begin{longtable}{p{0.28\linewidth}p{0.60\linewidth}}
\toprule
\textbf{Artifact type} & \textbf{Use} \\
\midrule
\code{fits_product} & Derived FITS/FIT/FTS not declared as display-only. \\
\code{fits_visual} & Display-only visual FITS, such as a rendered image \code{RGB_CUBE}. \\
\code{metadata_json} & Structured JSON other than the primary summary or manifest. \\
\code{preview_pdf} & PDF used as a preview or preview export. \\
\bottomrule
\end{longtable}

Prohibited official names are \code{json}, \code{csv}, \code{png}, \code{pdf},
\code{markdown}, \code{fits}, \code{directory}, \code{folder}, \code{preview},
\code{report}, \code{manifest}, and \code{notebook}: they are too generic for
the app contract.

\section{App-facing errors and blocks}
Primary reference: \path{references/v1-9-error-contract.md}.

A block or failure must provide output that an interface can present directly:
\begin{longtable}{p{0.30\linewidth}p{0.58\linewidth}}
\toprule
\textbf{Field} & \textbf{Requirement} \\
\midrule
\code{status} / \code{app_status} & Normalizes to \code{BLOCKED_CONTROLADO} or \code{FAIL}. \\
\code{errors[].kind} & An allowed category, not opaque free text. \\
\code{errors[].message} & A short, actionable human-readable message. \\
\code{command} & Redacted argv/cwd for traceability. \\
\code{inputs} & Inspected inputs, or an explanation that none apply. \\
\code{typed_artifacts} & Present even when empty for early blocks. \\
\code{next_actions} & At least one structured action for a block/failure. \\
\code{original_modified} & \code{false}, or explained \code{unknown}. \\
\code{short_summary} & Compact summary within \code{app_hints}. \\
\bottomrule
\end{longtable}

\subsection{Allowed kinds}
\begin{longtable}{p{0.35\linewidth}p{0.53\linewidth}}
\toprule
\textbf{Kind} & \textbf{Case} \\
\midrule
\code{missing_input} & Missing path or input. \\
\code{output_conflict} & Output is a directory, has an invalid parent, or is not writable. \\
\code{corrupt_input} & Suffix appears valid, but content is corrupt or truncated. \\
\code{missing_optional_backend} & Optional backend is absent. \\
\code{invalid_argument} & Invalid CLI argument. \\
\code{unsupported_format} & Format/family is unsupported by the selected route. \\
\code{permission_denied} & System permissions prevent reading or writing. \\
\code{timeout} & Backend/subprocess time limit exceeded. \\
\code{dependency_error} & Python dependency or runtime is unavailable. \\
\code{execution_error} & Execution fails after initial validation. \\
\code{tool_error} & Compatible fallback for unclassified legacy messages. \\
\bottomrule
\end{longtable}

\section{App hints and next actions}
v1.9 closes Debt B through focused capability profiles, improving explanation
without changing the scientific result.
\begin{longtable}{p{0.31\linewidth}p{0.57\linewidth}}
\toprule
\textbf{Field} & \textbf{Use} \\
\midrule
\code{short_summary} & Brief sentence for lists, cards, or notifications. \\
\code{severity} & \code{ok}, \code{warning}, \code{blocked}, or \code{error}. \\
\code{preview_artifact_types} & Previews the app can highlight. \\
\code{tags} & Route labels such as table, container, astronomy, expert-mode, or QA. \\
\code{next_actions} & Human next step: review columns, correct input, choose a backend, or read a manifest. \\
\bottomrule
\end{longtable}

\path{references/v1-9-app-hints-next-actions.md} records 17 closed targets,
including:
\begin{itemize}[leftmargin=1.5em]
\item General tables: \code{profile_table.py}.
\item General folders: \code{cross_domain_data_workbench.py}.
\item Time series: \code{timeseries_forecasting_workbench.py}.
\item QA/measurement: \code{physical_qa.py} and \code{photometry_noise_budget.py}.
\item Received notebooks: \code{coursework_notebook_fidelity_check.py}.
\item Expert astronomy: \code{catalog_workbench.py crossmatch-sky}.
\end{itemize}

\section{Capability exposure}
v1.9 freezes exposure modes to prevent narrow routes from appearing as ordinary
actions.
\begin{longtable}{p{0.26\linewidth}p{0.62\linewidth}}
\toprule
\textbf{Mode} & \textbf{Use} \\
\midrule
\code{normal} & Ordinary tables, documents, notebooks, containers, time series, or general QA. \\
\code{expert} & Advanced/domain-specific route requiring technical context. \\
\code{optional} & Optional backend or specific platform; not a core dependency. \\
\code{legacy} & IRAF, fxcor, iSTARMOD, and legacy coursework. \\
\code{maintainer} & Release gates, registry, smoke, synchronization, and probes. \\
\code{no_exponer} & Not a normal action; direct CLI or documented context only. \\
\bottomrule
\end{longtable}

For general inputs, \code{scientific_workflow_router.py} recommends normal
routes and demotes or rejects sensitive routes with
\code{normal_user_action=false}. The 22 sensitive routes in
\path{references/v1-9-exposure-modes.md} do not become normal v1.9 actions.

\section{v1.9 regressions}
\begin{longtable}{p{0.49\linewidth}p{0.39\linewidth}}
\toprule
\textbf{Gate} & \textbf{Coverage} \\
\midrule
\path{audit_v1_9_phase0_inventory_regression.py} & Charter, 59 rows, P1/P2 counts, decisions, and families. \\
\path{audit_v1_9_artifact_types_regression.py} & Official types, aliases, prohibited names, and Debt A targets. \\
\path{audit_v1_9_error_contract_regression.py} & Clean failures in \code{env_doctor.py}, \code{iwork_workbench.py}, and \code{duckdb_workbench.py}. \\
\path{audit_v1_9_dedupe_regression.py} & Internal helpers and reduced app-facing duplication. \\
\path{audit_v1_9_app_hints_regression.py} & 17 targets with hints and next actions. \\
\path{audit_v1_9_general_p2_regression.py} & General/adaptable routes with app-facing P2/P1 work. \\
\path{audit_v1_9_exposure_modes_regression.py} & Router and optional, expert, legacy, and maintainer modes. \\
\path{audit_v1_9_v2_0_deferrals_regression.py} & Eight protected v2.0 rows. \\
\path{audit_v1_9_more_app_like_regression.py} & Additional app-like regression coverage. \\
\path{audit_v1_9_app_like_extended_regression.py} & Realistic, ambiguous, and broken Phase 6 cases. \\
\path{audit_v1_9_docs_regression.py} & References, TeX/PDF guide, extracted text, and rendered pages. \\
\bottomrule
\end{longtable}

\section{Expanded app-like coverage}
Phase 6 expands coverage without real private data. Minimum families:
\begin{longtable}{p{0.34\linewidth}p{0.54\linewidth}}
\toprule
\textbf{Family} & \textbf{Expected behavior} \\
\midrule
Professional table & Useful parseable profiling with unchanged originals. \\
Administrative document & Intake with explicit warnings for weakly supported formats. \\
Inherited notebook & Preflight blocks or warns about \code{input()} and side effects. \\
Corrupt container & Warning or clean block, not false success. \\
Time series & Warnings for ambiguous dates, duplicates, or uncertain frequency. \\
Sensor/measurement & QA for suspicious ranges, NaN/Inf, or questionable units. \\
Expert astronomy route & Accurate warning/rejection for domain-specific data. \\
Absent optional backend & \code{BLOCKED_CONTROLADO}, not core \code{FAIL}. \\
Output conflict & Clean block with \code{output_conflict}. \\
\bottomrule
\end{longtable}
Recorded main outputs:
\begin{lstlisting}[style=promptstyle]
tmp/v1_9_debt_R_app_like_extended/app_like_extended_regression_summary.json
tmp/v1_9_debt_R_app_like_extended/app_like_extended_regression_report.md
\end{lstlisting}

\section{Protected v2.0 scope}
Reference: \path{references/v1-9-v2-0-deferrals.md}.

Reserved for v2.0:
\begin{itemize}[leftmargin=1.5em]
\item Major Scientific Workbench synchronization.
\item Complete UI-driven workflows and a formal app planner/skill execution engine.
\item Job history, previews, and statuses consumed directly by the app.
\item Joint skill/app documentation.
\item Interactive column, unit, style, and match selectors.
\item Visual confirmations and diffs before editing.
\item Complex optional panels for external backends.
\end{itemize}
Protected capability rows:
\begin{itemize}[leftmargin=1.5em]
\item \code{fits_rgb_batch.py}
\item \code{stilts_workbench.py}
\item \code{radial_velocity_workbench.py inspect}
\item \code{photometric_solution.py}
\item \code{apt_workbench.py}
\item \code{office_roundtrip.py docx-style-inventory}
\item \code{office_roundtrip.py docx-styled-replace}
\item \code{keynote_export.py}
\end{itemize}

\section{Testing a v1.9 capability}
Recommended pattern:
\begin{enumerate}[leftmargin=1.7em]
\item Create or copy anonymized inputs into a fresh \code{tmp/} run.
\item Run the router in dry-run mode when input is ambiguous.
\item Invoke the capability with \code{--summary-json}, \code{--manifest-json}, or \code{--run-dir} when available.
\item Validate JSON parsing and accurate \code{status} / \code{app_status}.
\item Check \code{typed_artifacts} against the frozen list.
\item Inspect \code{app_hints}, \code{next_actions}, \code{warnings}, and \code{errors}.
\item Confirm \code{original_modified=false}, or explained \code{unknown}.
\item Classify as \code{PASS}, \code{WARNING}, \code{BLOCKED_CONTROLADO}, \code{FAIL}, or audit \code{ROTO}.
\end{enumerate}
Examples:
\begin{lstlisting}[style=promptstyle]
python scripts/audit_v1_9_artifact_types_regression.py
python scripts/audit_v1_9_error_contract_regression.py
python scripts/audit_v1_9_app_like_extended_regression.py
\end{lstlisting}

\section{Historical Phase 7 closure}
Phase 7 requires:
\begin{itemize}[leftmargin=1.5em]
\item Expected v1.9 references with the required terms.
\item This guide in TeX and PDF.
\item Successful PDF compilation, extractable text containing key terms, and rendered representative pages or all pages.
\item Passing \code{scripts/audit_v1_9_docs_regression.py}.
\item No applicable drift from \code{scripts/sync_public_surface_docs.py --check}.
\item Scientific Workbench unchanged and the installed skill not synchronized.
\end{itemize}

\section{Toward Phases 8 and 9}
Phase 8 is the editable-tree release candidate: repeat v1.9 gates, affected
v1.8/v1.7/v1.6 regressions, and core smoke. Only Phase 9 may synchronize the
installed skill, after a release candidate without \code{FAIL} or \code{ROTO}.

This guide records the v1.9 documentation state, not final installation:
\begin{center}
\textbf{v1.9 consolidates the backend; v2.0 will synchronize Scientific Workbench.}
\end{center}
\end{document}
